\chapter[Referencial Teórico]{Referencial Teórico}
\label{cap:referencial}

Este capítulo reúne os fundamentos necessários à compreensão do trabalho:
o processamento de linguagem natural aplicado a textos clínicos, a tarefa de
extração de relações, as abordagens baseadas em regras e em aprendizado, e os
modelos de linguagem baseados em \foreignlanguage{english}{transformers}.

\section{Processamento de linguagem natural clínico}

O PLN clínico ocupa-se da análise automática de textos produzidos no cuidado em
saúde, como notas de evolução, sumários de alta e laudos. Esses textos
apresentam características que os distinguem da linguagem geral: vocabulário
técnico, abreviações não padronizadas, telegrafismo (omissão de elementos
gramaticais) e prevalência de construções negativas e de incerteza. Tais
propriedades tornam a transferência direta de ferramentas treinadas em domínio
aberto pouco eficaz, motivando recursos e modelos específicos do domínio.

\section{Extração de entidades e de relações}

A extração de informação em texto clínico costuma ser decomposta em duas
tarefas encadeadas. A primeira, \textbf{reconhecimento de entidades nomeadas}
(\foreignlanguage{english}{Named Entity Recognition}, NER), delimita trechos do
texto que mencionam conceitos de interesse e lhes atribui um tipo semântico. A
segunda, \textbf{extração de relações} (RE), determina se um par de entidades
mantém uma relação e de que tipo. Formalmente, dado um par ordenado de entidades
$(e_1, e_2)$ em um documento, a RE pode ser modelada como um problema de
classificação que associa ao par um rótulo dentro de um conjunto pré-definido.

Neste trabalho o espaço de rótulos é \texttt{negation\_of},
\texttt{associated\_with} e \texttt{no\_relation}. O rótulo
\texttt{no\_relation} é fundamental: como a anotação de corpora registra apenas
relações positivas, todo par de entidades não anotado deve ser tratado como
ausência de relação, o que introduz um forte desbalanceamento de classes.

\section{Abordagens baseadas em regras e a negação clínica}

Abordagens baseadas em regras codificam conhecimento linguístico explícito. Em
detecção de negação clínica, o algoritmo \textbf{NegEx} \cite{chapman_2001_negex}
tornou-se referência: a partir de um léxico de gatilhos de negação e de janelas
de contexto, classifica achados como negados ou afirmados. Sua simplicidade e
eficácia derivam do fato de a negação clínica ser altamente lexicalizada. Em
português, gatilhos como ``não'', ``sem'', ``nega'' e ``ausência de'', além de
adjetivos que embutem negação morfologicamente (``afebril'', ``anictérico''),
sustentam estratégias análogas (Apêndice~\ref{ap:lexico}). A alta lexicalização
sugere que o sinal de negação possa estar acessível mesmo a modelos não
especializados no domínio clínico, hipótese que este trabalho testa
diretamente.

\section{Abordagens baseadas em aprendizado de máquina}

Métodos aprendidos estimam, a partir de exemplos rotulados, a associação entre
características do par de entidades e o rótulo da relação. Modelos lineares como a
regressão logística, alimentados por características léxicas (por exemplo,
representações \foreignlanguage{english}{TF-IDF} do texto entre as entidades) e
estruturais (tipos das entidades, distância, direção), oferecem um compromisso
entre interpretabilidade e desempenho e são frequentemente empregados como
\emph{sanity check} antes de modelos mais custosos. Sua principal limitação é a
representação de palavras como unidades independentes, incapaz de capturar
dependências de longo alcance e contexto, limitação que motivou a transição da
área para modelos contextuais.

\section{Modelos baseados em \emph{transformers}}

A arquitetura \foreignlanguage{english}{Transformer} \cite{vaswani_2017_attention},
baseada em mecanismos de atenção, viabilizou modelos de linguagem pré-treinados
que produzem representações contextuais das palavras. O \textbf{BERT}
\cite{devlin_2019_bert} popularizou o paradigma de pré-treinamento seguido de
ajuste fino (\foreignlanguage{english}{fine-tuning}) para tarefas específicas.
Para o português, modelos multilíngues como o \textbf{XLM-R}
\cite{conneau_2020_xlmr} e, sobretudo, modelos especializados no domínio clínico
como o \textbf{BioBERTpt} \cite{schneider_2020_biobertpt}, adaptado a narrativas
clínicas e textos biomédicos em português brasileiro, constituem candidatos
naturais para a extração de relações clínicas, e é sobre eles que se apoia a
proposta deste trabalho.
